\chapter{Mã nguồn lõi: Count-Min Sketch}
\label{app:cms}

Count-Min Sketch là cấu trúc đếm tần suất xấp xỉ dùng trong pipeline để phát hiện quét
cổng (kho \texttt{K8sSecDaP-pipeline}, \texttt{libdsa/src/frequency/count\_min\_sketch.cpp}).
Cấu trúc gồm một bảng $d \times w$ ô đếm và $d$ hàm băm độc lập; \texttt{record} tăng một
ô trên mỗi hàng, còn \texttt{estimate} trả về \emph{giá trị nhỏ nhất} trong $d$ ô tương
ứng~--~đảm bảo ước lượng không bao giờ thấp hơn tần suất thật. Cơ sở lý thuyết về sai số
$\varepsilon$--$\delta$ đã trình bày ở Chương~\ref{ch:foundations}.

\begin{lstlisting}[language=C++,caption={Lõi \texttt{record}/\texttt{estimate}/\texttt{reset} của Count-Min Sketch.}]
// table_ : d hang x w cot (uint64_t); seeds_ : d hat giong bam doc lap.
void CountMinSketch::record(const void* key, size_t key_len) {
    for (size_t i = 0; i < depth_; ++i) {
        uint32_t col = hash(key, key_len, seeds_[i]) % width_;
        table_[i][col]++;                       // tang 1 o tren moi hang
    }
}

uint64_t CountMinSketch::estimate(const void* key, size_t key_len) const {
    uint64_t min_val = std::numeric_limits<uint64_t>::max();
    for (size_t i = 0; i < depth_; ++i) {
        uint32_t col = hash(key, key_len, seeds_[i]) % width_;
        min_val = std::min(min_val, table_[i][col]);  // lay MIN -> chan tren nhieu
    }
    return min_val;
}

void CountMinSketch::reset() {                  // dat lai khi het cua so tumbling
    for (auto& row : table_) std::fill(row.begin(), row.end(), 0);
}

size_t CountMinSketch::memory_usage() const {   // bo nho khong doi theo so khoa
    return depth_ * width_ * sizeof(uint64_t) + seeds_.size() * sizeof(uint32_t);
}
\end{lstlisting}

Với cấu hình của đồ án ($d = 5$, $w = 2048$), bộ nhớ là $5 \times 2048 \times 8 \approx
80$~KB và \emph{không đổi} dù số địa chỉ nguồn phân biệt tăng lên~--~đây chính là ưu thế
của sketch so với đếm chính xác bằng \texttt{HashMap} (bộ nhớ tuyến tính theo số khóa).
